% Five-page chapter. Source reconciliation and layout notes:
% output/ch9/five_pages_20260925/CH9_revision_notes.md
\chapter{Systematic Uncertainties}
\label{chapter:systematics_uncertainties}

Systematic uncertainties in this analysis includes the experimental and theoretical modelling uncertainties.
Their sources and evaluation are discussed in this chapter.

\section{Experimental uncertainties}
\label{sec:ch9_experimental}

Experimental uncertainties originate from the detector response and object calibrations. They are propagated through the event weights or reconstructed objects.
The evaluation and treatment of these uncertainties are summarised below.

\subsubtitle{Luminosity and pile-up}
\label{subsec:ch9_lumi_pileup}
The 2015-2018 integrated luminosity has an uncertainty of $0.83\%$, as discussed in Section~\ref{sec:dataset_selection}.
It is applied to the signal and backgrounds normalisation.
The pile-up uncertainty is evaluated by varying the data-derived $\mu$ profile used to calculate the pile-up correction weights for simulated events. Its nominal rescaling factor, $1/1.03$, is changed to $1/1.07$ and $1/0.99$.

\subsubtitle{Hadronic tau energy scale and efficiencies}
\label{subsec:ch9_tau}
The TES uncertainty in the low-$\pT$ region is measured using $\ztautau$ events, while the uncertainty in the high-$\pT$ region is evaluated using test-beam data and MC comparisons~\cite{ATLAS-CONF-2017-029,ATL-PHYS-PUB-2015-045}.
Efficiency uncertainties are evaluated for $\hadtau$ reconstruction and RNN identification using $\ttbar$ events, as well as for electron rejection.
For the \texttt{Tight} ID working point, the scale-factor uncertainties are $5$--$11\%$.

\subsubtitle{Jets and missing transverse momentum}
\label{subsec:ch9_jets_met}
The JES and JER uncertainties are evaluated using the calibration methods described in Section~\ref{sec:ch6_jets}.
The $\mpt$ uncertainty is evaluated by recalculating the $\mpt$ for each uncertainty from the hard-object momenta and the soft term.
The soft-term uncertainties are derived from measurements in $\zee$ events~\cite{ATLAS2025METPerformance}.

\subsubtitle{Jet flavour tagging}
\label{subsec:ch9_ftag}
The flavour-tagging uncertainties are evaluated separately for the $b$-jet tagging efficiency and the $c$-jet and light-jet mistagging efficiencies.
They include calibration sample statistics, background contamination and modelling uncertainties in the efficiency measurements described in Section~\ref{sec:ch6_btag}.
% An eigenvector decomposition of the calibration covariance matrix provides independent variations that preserve correlations between bins.
% The \texttt{Medium} reduction scheme is used.
% The resulting scale-factor variations affect both $\bTagSR$ and $\bVetoSR$ through the probabilities for jets to pass or fail the tagging requirement.

\subsubtitle{Electrons and muons}
\label{subsec:ch9_leptons}
Electron efficiency uncertainties are evaluated using $\zee$ and $J/\psi\to ee$ events, while energy-scale and resolution uncertainties are obtained by comparing the measured and simulated calorimeter response~\cite{ATLAS2019ElectronPerformance}.
The muon efficiency, momentum-scale and resolution uncertainties are determined using $\zmumu$ and $J/\psi\to\mu\mu$ events~\cite{ATLAS2021MuonID,ATLAS2023MuonCalibration}.
In particular, the resolution uncertainty has a sizeable effect on the high-$\pT$ muons in this analysis. 



\section{Theoretical uncertainties}
\label{sec:ch9_theory}

Considering that the background fit uses the total event yield in each region, the transfer-factor (TF) method is introduced to evaluate the modelling uncertainties in the $\wjets$ and top-quark backgrounds.
Given that the $\wjets$ (top-quark background) contamination in $\CRTop$ ($\CRW$) is below $5\%$ ($2\%$), the modelling uncertainties on these small contributions are neglected.
The uncertainties types for the $\wjets$ and top-quark backgrounds are then introduced.

\subsection{Transfer-factor method}
\label{subsec:ch9_transfer}

The background yield in an SR or VR is estimated as
\begin{align}
    N_{bg}^{SR(VR),data}
    &=\lambda\times N_{bg}^{SR(VR),MC}
    =TF\times N_{bg}^{CR,data}.
    \label{eq:ch9_transfer_prediction}
\end{align}
The NF $\lambda$ is introduced as a free-floating parameter as described in Equation~\ref{eq:ch8_lambda}. $TF$ is defined by
\begin{align}
    TF=\frac{N_{bg}^{SR(VR),MC}}{N_{bg}^{CR,MC}}.
    \label{eq:ch9_transfer}
\end{align}

Suppose a theory uncertainty changes the background yield by $\Delta N_{bg}=N_{bg}^{sysMC}-N_{bg}^{nomMC}$, where $nomMC$ and $sysMC$ denote the nominal and alternative MC predictions, respectively.
Propagating this change to $TF$ gives
\begin{align}
    TF+\Delta TF
    &=\frac{N_{bg}^{SR(VR),MC}+\Delta N_{bg}^{SR(VR),MC}}
    {N_{bg}^{CR,MC}+\Delta N_{bg}^{CR,MC}}\nonumber\\
    &=TF\times\frac{N_{bg}^{SR(VR),sysMC}}{N_{bg}^{SR(VR),nomMC}}
    \bigg/\frac{N_{bg}^{CR,sysMC}}{N_{bg}^{CR,nomMC}}.
    \label{eq:ch9_transfer_variation}
\end{align}
If an uncertainty produces the same fractional change in both regions, it cancels in the double ratio, giving $\Delta TF=0$.
The remaining relative uncertainty is applied only to the SR or VR yield through the factor $1\pm\Delta TF/TF$.
% For alternative samples with low statistics, smoothing and relaxed jet-multiplicity requirements are used where necessary to stabilise the uncertainty estimates.

\subsection{$\wjets$}
\label{subsec:ch9_wjets}
\begin{itemize}[nosep,leftmargin=*]
    \item \textbf{QCD scales:} the uncertainty is evaluted from the envelope of seven-point predictions obtained by varying the matrix-element and parton-shower renormalisation and factorisation scales~\cite{PMGR-2021-01}.
    \item \textbf{PDF:} the uncertainty is evaluated using the Hessian NNPDF3.0 NNLO uncertainty scheme, including alternative PDF sets and two alternative $\alpha_s$ values.
    % \item \textbf{Electroweak corrections:} the uncertainty is given by the approximate NLO electroweak virtual corrections.
    \item \textbf{$W+\mathrm{HF}/W+\mathrm{LF}$ ratio:} an uncertainty is assigned to heavy-to-light-flavour ratio when extrapolating from the $b$-vetoed $\CRW$ to regions with $b$ jets.
    A value of $30\%$ is adopted, motivated by flavour-dependent normalisation differences observed in an earlier ATLAS measurement~\cite{HIGG-2020-20}.
\end{itemize}

\subsection{Top-quark background}
\label{subsec:ch9_top}
\begin{itemize}[nosep,leftmargin=*]
    \item \textbf{QCD scales, PDF and PDF $\alpha_s$:} the uncertainties are evaluated by reweighting the nominal samples, as for $\wjets$.
    \item \textbf{Matrix element (NLO matching):} the uncertainty is evaluated by changing the definition of the region in which \textsc{Pythia} shower emissions are vetoed to avoid overlap with radiation generated by \textsc{Powheg}.
    \item \textbf{Parton shower and hadronisation:} the uncertainty is obtained by comparing the nominal samples with alternative samples using \textsc{Herwig}~7.
    \item \textbf{ISR (FSR) $\alpha_s$:} the uncertainty is evaluated using internal weights for initial-state (final-state) radiation.
    \item \textbf{$h_{\mathrm{damp}}$:} the uncertainty is evaluated by varying $h_{\mathrm{damp}}$, which regulates hard radiation recoiling against the $\ttbar$ system~\cite{ATL-PHYS-PUB-2020-023}.
    \item \textbf{$\ttbar$ NNLO:} the uncertainty is derived by reweighting the nominal sample to the NNLO prediction~\cite{Czakon:2017wor}.
    \item \textbf{$\ttbar$ shower recoil:} the uncertainty is evaluated by assigning the recoil from subsequent gluon emissions to the parent top quark instead of the $b$ quark in the top quark decay.
    \item \textbf{$Wt$--$\ttbar$ interference:} the uncertainty is evaluated by comparing the nominal $Wt$ sample with an alternative sample in which the overlapping $\ttbar$ contribution is subtracted~\cite{Frixione:2008yi}.
\end{itemize}

\subsection{Other backgrounds and signal}
\label{subsec:ch9_signal}
The $\znunu$ background with a jet misidentified as a $\hadtau$ is assigned with a $100\%$ normalisation uncertainty, while diboson and the remaining $\ztautau$ backgrounds are assigned with a conservative $30\%$ normalisation uncertainty.
Signal uncertainties are evaluated using QCD scale and PDF variations, resulting in a $10$--$20\%$ uncertainty.
The background theory uncertainties are summarised in Tables~\ref{tab:ch9_theory_res} and~\ref{tab:ch9_theory_nonres} for the resonant and non-resonant regions, respectively.

\clearpage
% The two numerical tables are inserted below from the checked INT source.
\begin{table}[!htbp]
    \centering
    \begin{spacing}{1}
    \captionsetup{font=footnotesize}
    \setlength{\tabcolsep}{3pt}
    \renewcommand{\arraystretch}{0.95}
    \fontsize{9}{10.5}\selectfont
    \caption[Theory uncertainties in the resonant regions]{Theory uncertainties considered in the resonant regions. The last four sources are applied to the overall normalisation rather than the TF.}
    \label{tab:ch9_theory_res}
    \label{tab:ch9_theory_transfer}
    \begin{tabular*}{\textwidth}{@{\extracolsep{\fill}}lccc@{}}
        \toprule
        Source & \multicolumn{3}{c}{$\Delta TF/TF$ Up (Down)} \\
        & $\bTagSR$ & $\VRW$ & $\VRTop$ \\
        \midrule
        $\wjets$ QCD scales & $0.10 (-0.17)$ & $0.04 (-0.04)$ & --- \\
        $\wjets$ PDF & $<0.001$ & $0.007$ & --- \\
        $\wjets$ PDF $\alpha_s$ & $-0.02 (-0.001)$ & $0.02 (0.001)$ & --- \\
        $\wjets$ EW correction & $0.04$ & $-0.02$ & --- \\
        \midrule
        Top NLO matching & $-0.06$ & $0.08$ & $0.23$ \\
        Top parton shower & $0.16$ & $0.1$ & $0.11$ \\
        Top $\mu_F$ & $0.12 (-0.09)$ & $0.03 (-0.02)$ & $0.01 (-0.01)$ \\
        Top $\mu_R$ & $0.15 (-0.19)$ & $-0.01 (0.01)$ & $0.005 (-0.008)$ \\
        Top ISR & $-0.04 (-0.01)$ & $0.01 (-0.01)$ & $-0.01 (0.01)$ \\
        Top FSR & $0.19 (-0.24)$ & $-0.04 (0.08)$ & $-0.02 (0.22)$ \\
        Top PDF & $-0.051$ & $0.043$ & $0.032$ \\
        Top PDF $\alpha_s$ & $0.01 (-0.01)$ & $0.01 (-0.004)$ & $0.002 (0.0002)$ \\
        Top $h_{\mathrm{damp}}$ & $0.09$ & $0.20$ & $0.15$ \\
        $\ttbar$ NNLO & $-0.05$ & $0.02$ & $0.01$ \\
        $\ttbar$ shower recoil & $-0.08$ & $0.04$ & $-0.04$ \\
        $Wt$ interference & $-0.1$ & $-0.18$ & $-0.25$ \\
        \midrule
        $\znunu$ with fake $\hadtau$ & $1.0$ & $0$ & $1.0$ \\
        Diboson modelling & $0.3$ & $0.3$ & $0.3$ \\
        $\zjets$ modelling & $0.3$ & $0.3$ & $0.3$ \\
        $\wjets$ heavy flavour\textsuperscript{*} & $0.3$ & $0.3$ & $0.3$ \\
        \bottomrule
    \end{tabular*}
    \vspace{1.6em}
    \caption[Theory uncertainties in the non-resonant regions]{Theory uncertainties considered in the non-resonant regions. The last four sources are applied to the overall normalisation rather than the TF.}
    \label{tab:ch9_theory_nonres}
    \begin{tabular*}{\textwidth}{@{\extracolsep{\fill}}lccc@{}}
        \toprule
        Source & \multicolumn{3}{c}{$\Delta TF/TF$ Up (Down)} \\
        & $\bTagSR$ & $\VRW$ & $\VRTop$ \\
        \midrule
        $\wjets$ QCD scales & $0.14 (-0.16)$ & $0.01 (-0.03)$ & --- \\
        $\wjets$ PDF & $-0.003$ & $0.002$ & --- \\
        $\wjets$ PDF $\alpha_s$ & $0.03 (-0.03)$ & $0.02 (-0.004)$ & --- \\
        $\wjets$ EW correction & $0.01$ & $0.01$ & --- \\
        \midrule
        Top NLO matching & $-0.15$ & $0.01$ & $-0.03$ \\
        Top parton shower & $-0.01$ & $-0.14$ & $-0.07$ \\
        Top $\mu_F$ & $-0.05 (0.05)$ & $0.01 (-0.01)$ & $0.003 (-0.004)$ \\
        Top $\mu_R$ & $-0.01 (0.02)$ & $0.04 (-0.05)$ & $0.01 (-0.02)$ \\
        Top ISR & $0.16 (-0.13)$ & $-0.0003 (0.002)$ & $-0.01 (<0.001)$ \\
        Top FSR & $0.01 (-0.25)$ & $0.06 (-0.14)$ & $0.003 (-0.04)$ \\
        Top PDF & $-0.025$ & $-0.002$ & $0.003$ \\
        Top PDF $\alpha_s$ & $-0.01 (0.002)$ & $0.005 (-0.001)$ & $0.001 (<0.001)$ \\
        Top $h_{\mathrm{damp}}$ & $-0.07$ & $-0.01$ & $0.05$ \\
        $\ttbar$ NNLO & $-0.003$ & $-0.004$ & $-0.01$ \\
        $\ttbar$ shower recoil & $-0.13$ & $0.09$ & $0.03$ \\
        $Wt$ interference & $0.03$ & $-0.04$ & $-0.009$ \\
        \midrule
        $\znunu$ with fake $\hadtau$ & $1.0$ & $0$ & $1.0$ \\
        Diboson modelling & $0.3$ & $0.3$ & $0.3$ \\
        $\zjets$ modelling & $0.3$ & $0.3$ & $0.3$ \\
        $\wjets$ heavy flavour\textsuperscript{*} & $0.3$ & $0.3$ & $0.3$ \\
        \bottomrule
    \end{tabular*}
    \par\smallskip
    \raggedright
    \textsuperscript{*} The $\wjets$ heavy-flavour uncertainty is applied independently in the resonant and non-resonant regions.
    \end{spacing}
\end{table}
